\documentclass[aspectratio=169]{beamer}
\usepackage{etex} % fixes new-dimension error
\usepackage{lmodern}
\usepackage[T1]{fontenc}

\usepackage{graphicx,amsmath}
\usepackage{stmaryrd} % cf. interleave
\usepackage{booktabs}
\usepackage{amscd}
\usepackage{multicol}
\usepackage[absolute,overlay]{textpos}
\usepackage{alltt}
\usepackage{proof}
\usepackage{subcaption}
\usepackage{tikz}
\usepackage{tikz-cd}
\usepackage[new]{old-arrows}
\usepackage[all]{xy}
\usepackage{pgfplots}
\usepackage{textcomp}

\usepackage{transparent}
\usepackage{xspace}
\usepackage{listings}
\usepackage{pdfpages}
\usepackage{relsize}

%%%%%%%%%%%%% Macros
\newcommand{\Ban}{\catfont{Ban}}
\newcommand{\Met}{\catfont{Met}}
\newcommand{\Shuff}{\mathrm{Sf}}
\newcommand{\Cats}{\catfont{Cat}}
\newcommand{\VCat}{\mathcal{V}\text{-}\Cats}
\newcommand{\VCatSy}{\mathcal{V}\text{-}\Cats_{\mathsf{sym}}}
\newcommand{\VCatSe}{\mathcal{V}\text{-}\Cats_{\mathsf{sep}}}
\newcommand{\VCatSS}{\mathcal{V}\text{-}\Cats_{\mathsf{sym,sep}}}
%%%% Categories
\newcommand{\catfont}[1]{\mathsf{#1}}
\newcommand{\cop}{\catfont{op}}
\newcommand{\Law}{\catfont{Law}}
\newcommand{\catV}{\catfont{V}}
\newcommand{\catX}{\catfont{X}}
\newcommand{\catC}{\catfont{C}}
\newcommand{\catD}{\catfont{D}}
\newcommand{\catA}{\catfont{A}}
\newcommand{\catB}{\catfont{B}}
\newcommand{\catI}{\catfont{I}}
\newcommand{\Set}{\catfont{Set}}
\newcommand{\Top}{\catfont{Top}}
\newcommand{\Pos}{\catfont{Pos}}
\newcommand{\Inj}{\catfont{Inj}}
\newcommand{\Det}{\catfont{RMhat}}
\newcommand{\CoAlg}[1]{\catfont{CoAlg}\left (#1 \right )}
\newcommand{\Mon}{\catfont{Mon}}
\newcommand{\Mnd}{\catfont{Mnd}(\catC)}
\newcommand{\SMnd}{\catfont{Mnd}(\Set)}
\newcommand{\CLat}{\catfont{CLat}}
\newcommand{\Stone}{\catfont{Stone}}
\newcommand{\Spectral}{\catfont{Spectral}}
\newcommand{\CompHaus}{\catfont{CompHaus}}
\newcommand{\Subs}[2]{\catfont{Sub}_{}}
\newcommand{\Cone}{\catfont{Cone}}
\newcommand{\StComp}{\catfont{StablyComp}}
\newcommand{\PosC}{\catfont{PosComp}}
\newcommand{\Haus}{\catfont{Haus}}
\newcommand{\Meas}{\catfont{Meas}}
\newcommand{\Ord}{\catfont{Ord}}
\newcommand{\EndoC}{[\catC,\catC]}
%% General functors
\newcommand{\funfont}[1]{#1}
\newcommand{\funF}{\funfont{F}}
\newcommand{\funU}{\funfont{U}}
\newcommand{\funG}{\funfont{G}}
\newcommand{\funT}{\funfont{T}}
\newcommand{\funI}{\funfont{I}}
%% Particular kinds of functors
\newcommand{\sfunfont}[1]{\mathrm{#1}}
\newcommand{\Pow}{\sfunfont{P}}
\newcommand{\Dist}{\sfunfont{D}}
\newcommand{\Maybe}{\sfunfont{M}}
\newcommand{\List}{\sfunfont{L}}
\newcommand{\UForg}{\sfunfont{U}}
\newcommand{\Forg}[1]{\sfunfont{U}_{#1}}
\newcommand{\Id}{\sfunfont{Id}}
\newcommand{\Vie}{\sfunfont{V}}
\newcommand{\Disc}{\funfont{D}}
\newcommand{\Weight}{\sfunfont{W}}
\newcommand{\homf}{\sfunfont{hom}}
\newcommand{\Yoneda}{\sfunfont{Y}}
%% Diagram functors
\newcommand{\Diag}{\mathscr{D}}
\newcommand{\KDiag}{\mathscr{K}}
\newcommand{\LDiag}{\mathscr{L}}
%% Monads
\newcommand{\monadfont}[1]{\mathbb{#1}}
\newcommand{\monadT}{\monadfont{T}}
\newcommand{\monadS}{\monadfont{S}}
\newcommand{\monadU}{\monadfont{U}}
\newcommand{\monadH}{\monadfont{H}}
\newcommand{\str}{\mathrm{str}}
%% Adjunctions
\newcommand\adjunct[2]{\xymatrix@=8ex{\ar@{}[r]|{\top}\ar@<1mm>@/^2mm/[r]^{{#2}}
& \ar@<1mm>@/^2mm/[l]^{{#1}}}}
\newcommand\adjunctop[2]{\xymatrix@=8ex{\ar@{}[r]|{\bot}\ar@<1mm>@/^2mm/[r]^{{#2}}
& \ar@<1mm>@/^2mm/[l]^{{#1}}}}
%% Retractions
\newcommand\retract[2]{\xymatrix@=8ex{\ar@{}[r]|{}\ar@<1mm>@/^2mm/@{^{(}->}[r]^{{#2}}
& \ar@<1mm>@/^2mm/@{->>}[l]^{{#1}}}}
%% Limits
\newcommand{\pv}[2]{\langle #1, #2 \rangle}
\newcommand{\limt}{\mathrm{lim}}
\newcommand{\pullbackcorner}[1][dr]{\save*!/#1+1.2pc/#1:(1,-1)@^{|-}\restore}
\newcommand{\pushoutcorner}[1][dr]{\save*!/#1-1.2pc/#1:(-1,1)@^{|-}\restore}
%% Colimits
\newcommand{\colim}{\mathrm{colim}}
\newcommand{\inl}{\mathrm{inl}}
\newcommand{\inr}{\mathrm{inr}}
%% Distributive categories
\newcommand{\distr}{\mathrm{dist}}
\newcommand{\undistr}{\mathrm{undist}}
%% Closedness
\newcommand{\curry}[1]{\mathrm{curry}{#1}}
\newcommand{\app}{\mathrm{app}}
%% Misc. operations
\newcommand{\const}[1]{\underline{#1}}
\newcommand{\comp}{\cdot}
\newcommand{\id}{\mathrm{id}}
%% Factorisations
\newcommand{\EClass}{E}
\newcommand{\MClass}{M}
\newcommand{\MConeClass}{\mathcal{M}}
%%%%%%%%%%%%%%%% End of Categorical Stuff

%%%% Misc
%% Operations
\newcommand{\blank}{\, - \,}
\newcommand{\sem}[1]{\llbracket #1 \rrbracket}
\newcommand{\closure}[1]{\overline{#1}}
\DeclareMathOperator{\img}{\mathrm{im}}
\DeclareMathOperator{\dom}{\mathrm{dom}}
\DeclareMathOperator{\codom}{\mathrm{codom}}
%% Sets of numbers
\newcommand{\N}{\mathbb{N}}
\newcommand{\Z}{\mathbb{Z}}
\newcommand{\Nats}{\mathbb{N}}
\newcommand{\Reals}{\mathbb{R}}
\newcommand{\Rz}{\Reals_{\geq 0}}
\newcommand{\Complex}{\mathbb{C}}
%% Writing
\newcommand{\cf}{\emph{cf.}}
\newcommand{\ie}{\emph{i.e.}}
\newcommand{\eg}{\emph{e.g.}}
\newcommand{\df}[1]{\emph{\textbf{#1}}}
%%%%%%%%%%%%%%%% End of Misc

%%%% Programming Stuff
%% Types
\newcommand{\typefont}[1]{\mathbb{#1}}
\newcommand{\typeOne}{1}
\newcommand{\typeTwo}{2}
\newcommand{\typeA}{\typefont{A}}
\newcommand{\typeX}{\typefont{X}}
\newcommand{\typeB}{\typefont{B}}
\newcommand{\typeC}{\typefont{C}}
\newcommand{\typeV}{\typefont{V}}
\newcommand{\typeD}{\typefont{D}}
\newcommand{\typeI}{\typefont{I}}
%% RuleName
\newcommand{\rulename}[1]{(\mathrm{#1})}
%% Sequents
\newcommand{\jud}{\vdash}
\newcommand{\vljud}{\rhd}
\newcommand{\cojud}{\vdash_{\co}}
\newcommand{\vl}{\mathtt{v}}
\newcommand{\co}{\mathtt{c}}
% Program font
\newcommand{\prog}[1]{\mathtt{#1}}
\newcommand{\pseq}[3]{#1 \leftarrow #2; #3}
\newcommand{\ppm}[4]{(#1,#2) \leftarrow #3; #4}
\newcommand{\pinl}[1]{\prog{inl}(#1)}
\newcommand{\pinr}[1]{\prog{inr}(#1)}
\newcommand{\pcase}[4]{\prog{ case } #1 \prog{ of } \pinl{#2} \Rightarrow #3 ; \pinr{#2} \Rightarrow #4}
%% Sets of terms
\newcommand{\ValuesBP}[2]{\mathsf{Values}(#1, #2)}
\newcommand{\TermsBP}[2]{\mathsf{Terms}(#1, #2)}
\newcommand{\closValP}[1]{\ValuesBP{\emptyset}{#1}}
\newcommand{\closTermP}[1]{\TermsBP{\emptyset}{#1}}
\newcommand{\closVal}{\closValP{\typeA}}
\newcommand{\closTerm}{\closTermP{\typeA}}
%% Contextual equivalence
\newcommand{\ctxeq}{\equiv_{\prog{ctx}}}
%%%% End of Programming Stuff

%------ Setting lecture info ----------------------------------------------
\newcounter{lectureID}
\stepcounter{lectureID}
\newcommand{\getLecture}{\arabic{lectureID}\xspace}
\newcommand{\setLectureBasic}[1]{
  \title{
    #1
    }
  \author{Jos\'{e} Proen\c{c}a}
  \institute{CISTER -- U.Porto, Porto, Portugal
            \hfill 
            \begin{tabular}{r@{}}
            \url{https://cister-labs.github.io/alg2324}
            \end{tabular}
            }
  \date{Algorithms (CC4010) 2023/2024}
  % logos of institutions
  \titlegraphic{
    \begin{textblock*}{5cm}(6.0cm,6.80cm)
       \includegraphics[scale=0.18]{images/fcup}\hspace*{.85cm}~%
    \end{textblock*}
    % \begin{textblock*}{5cm}(9.4cm,7.45cm)
    %   \includegraphics[scale=0.50]{images/dcc}
    % \end{textblock*}
  }  
}
\newcommand{\setLecture}[2]{\setcounter{lectureID}{#1}\setLectureBasic{#1. #2}}

%------ Counters for exercises ----------------------------------------------
\newcounter{cExercise}
\newcommand{\exercise}{\stepcounter{cExercise}Ex.\,\arabic{lectureID}.\arabic{cExercise}:\xspace}
\newcommand{\exerciseBack}{\addtocounter{cExercise}{-1}}
\newcommand{\exerciseAdd}{\stepcounter{cExercise}}
\newcommand{\doExercise}[3][0mm]{\begin{exampleblock}{\exercise #2}\wrap{\rule{0pt}{#1}}#3\end{exampleblock}}
\newcommand{\doSimpleExercise}[2][0mm]{\begin{exampleblock}{}\wrap{\rule{0pt}{#1}}\structure{\textbf{\exercise} #2}\end{exampleblock}}

% Slide
\newenvironment{slide}[1]{\begin{frame}\frametitle{#1}}{\end{frame}}

% Misc by José
\newcommand{\wrap}[2][]{\begin{tabular}[#1]{@{}c@{}}#2\end{tabular}}
\newcommand{\mwrap}[1]{\ensuremath{\begin{array}{@{}c@{}}#1\end{array}}}
\def\trans#1{\xrightarrow{#1}}  % - a - > 
\def\Trans#1{\stackrel{#1}{\Longrightarrow}} % =a=> 
\newcommand{\transp}[2][35]{\color{fg!#1}#2}
\newcommand{\transpt}[2][.35]{\tikz{\node[inner sep=1pt,fill opacity=0.5]{#2}}}
\newcommand{\faded}[2][0.4]{{\transparent{#1}#2}} % alternative to "transp" using transparent package
\newcommand{\set}[1]{\left\{ #1 \right\}} % {a,b,...z}
\newcommand{\mi}[1]{\ensuremath{\mathit{#1}}\xspace}
\newcommand{\mf}[1]{\ensuremath{\mathsf{#1}}\xspace}
% \newcommand{\gold}[1]{\textcolor{darkgoldenrod}{#1}\xspace}


%------ using color ---------------------------------------------------------
\definecolor{goldenrod}{rgb}{.80392 .60784 .11373}
\definecolor{darkgoldenrod}{rgb}{.5451 .39608 .03137}
\definecolor{brown}{rgb}{.15 .15 .15}
\definecolor{darkolivegreen}{rgb}{.33333 .41961 .18431}
\definecolor{myGray}{gray}{0.85}
%
%
\newcommand{\red}[1]{\textcolor{red!80!black}{#1}\xspace}
\newcommand{\blue}[1]{\textcolor{blue}{#1}\xspace}
\newcommand{\gold}[1]{\textcolor{darkgoldenrod}{#1}\xspace}
\newcommand{\gray}[1]{\textcolor{myGray}{#1}\xspace}
% \def\alert#1{{\darkgoldenrod #1}}
% \def\alert#1{{\alert{#1}}}
%\def\brw#1{{\brown #1}}
% \def\structure#1{{\blue #1}}
% \def\tstructure#1{\textbf{\darkblue #1}}
%%\def\gre#1{{\green #1}}
\def\gre#1{{\darkolivegreen #1}}
\def\gry#1{{\textcolor{gray}{#1}}}
\def\rdb#1{{\red #1}}
\def\st{\mathbf{.}\,}
\def\laplace#1#2{*\txt{\mbox{ \fcolorbox{black}{myGray}{$\begin{array}{c}\mbox{#1}\\\\#2\\\\\end{array}$} }}}
%\newcommand{\galois}[2]{#1\; \dashv\; #2}



% ----- from LSB
\def\Act{N}
\def\cnil{\mathbf{0}}
\def\cpf#1#2{#1 . #2}                           % a.P
\def\cou#1#2{#1 \mathbin{+} #2}                 % P + Q
%\def\crt#1#2{\mathbin{#1 \setminus_{#2}}}       % P \ A
%\def\crtt#1#2{\mathbin{#1 \setminus\!\setminus_{#2}}}       % P \ A
%\def\crt#1#2{\mathsf{new}\, #2\;  #1}       % P \ A
\def\crt#1#2{#1 \backslash #2}       % P \ A
%\def\crn#1#2{\{#2\}\, #1}                  % P[f]
\def\ainv#1{\overline{#1}}
\def\rtran#1{\stackrel{#1}{\longrightarrow}}
% \def\pair#1{\const{#1}}
\def\pair#1{\langle #1 \rangle}
\def\asor{\mathbin{|}}                    % A | B
\def\setdef#1#2{\mathopen{\{} #1 \asor #2 \mathclose{\}}}
\def\imp{\mathbin{\Rightarrow}}
\def\dimp{\mathbin{\Leftrightarrow}}
\def\rimp{\mathbin{\Leftarrow}}
\def\rra{\longrightarrow}
\def\rcb#1#2#3#4{\def\nothing{}\def\range{#3}\mathopen{\langle}#1 \ #2 \ \ifx\range\nothing::\else: \ #3 :\fi \ #4\mathclose{\rangle}}
\def\aconv#1{#1^{\circ}} 
\def\abv{\stackrel{\rm abv}{=}}


\def\crn#1#2{\mathbin{#1[#2]}}                  % P[f]
\def\couit#1#2{\Sigma_{#1}#2}                  %  + i=1,n
\def\cpar#1#2{#1 \mid #2}                       %  |
\def\ctpar#1#2{#1 \parallel #2}                       %  |
\def\cpars#1#2#3{#1 \mid_{#3} #2}               %  |S

% Spliting frames in 2 columns
\newcommand{\frsplit}[3][.48]{
  \begin{columns}%[T] % align columns
  \begin{column}{#1\textwidth} #2 \end{column} ~~~
  \begin{column}{#1\textwidth} #3 \end{column} \end{columns}
}
\newcommand{\frsplitdiff}[4]{
  \begin{columns}%[T] % align columns
  \begin{column}{#1\textwidth} #3 \end{column} ~~~
  \begin{column}{#2\textwidth} #4 \end{column} \end{columns}
}
\newcommand{\frsplitt}[3][.48]{
  \begin{columns}[T] % align columns
  \begin{column}{#1\textwidth} #2 \end{column} ~~~
  \begin{column}{#1\textwidth} #3 \end{column} \end{columns}
}
\newcommand{\col}[2][.48]{\begin{column}{#1\textwidth} #2 \end{column}}
\newcommand{\colb}[3][.48]{\begin{column}{#1\textwidth} \begin{block}{#2} #3 \end{block} \end{column}}

% Spliting frames in 3 columns
\newcommand{\frsplitthree}[4][.31]{
  \begin{columns}%[T] % align columns
  \begin{column}{#1\textwidth} #2 \end{column} ~~~
  \begin{column}{#1\textwidth} #3 \end{column} ~~~
  \begin{column}{#1\textwidth} #4 \end{column} \end{columns}
}
\newcommand{\frsplitdiffthree}[5]{
  \begin{columns}%[T] % align columns
  \begin{column}{#1\textwidth} #3 \end{column} ~~~
  \begin{column}{#1\textwidth} #4 \end{column} ~~~
  \begin{column}{#2\textwidth} #5 \end{column} \end{columns}
}
\newcommand{\frsplittthree}[4][.32]{
  \begin{columns}[T] % align columns
  \begin{column}{#1\textwidth} #2 \end{column} ~
  \begin{column}{#1\textwidth} #3 \end{column} ~
  \begin{column}{#1\textwidth} #4 \end{column} \end{columns}
}


\newcommand{\typerule}[4][]{\ensuremath{\begin{array}[#1]{c}\textsf{\scriptsize ({#2})} \\#3 \\\hline\raisebox{-3pt}{\ensuremath{#4}}\end{array}}}
\newcommand{\styperule}[3][]{\ensuremath{\begin{array}[#1]{c} #2 \\[0.5mm]\hline\raisebox{-4pt}{\ensuremath{#3}}\end{array}}}
\newcommand{\shrk}{\vspace{-3mm}}

\def\caixa#1{\medskip
  \begin{center}
  \fbox{\begin{minipage}{0.9\textwidth}\protect{#1}\end{minipage}}
  \end{center}}

\newcommand{\mybox}[2][0.9]{
  \begin{minipage}{#1\textwidth}\begin{block}{}\centering #2\end{block}\end{minipage}}
\newcommand{\mycbox}[2][0.9]{
  {\\[-5mm]\centering\mybox[#1]{#2}\\[-5mm]}}

%%%%% Tikz
% \usetikzlibrary{arrows.meta, calc, fit, tikzmark}
\usetikzlibrary{%
  positioning
 ,patterns
 ,arrows
 ,arrows.meta
 ,automata
 ,calc
 ,shapes
 ,fit
 ,tikzmark
 ,fadings
 ,decorations.pathreplacing
 ,plotmarks
% ,pgfplots.groupplots
 ,decorations.markings
}
% \tikzset{shorten >=1pt,node distance=2cm,on grid,auto,initial text={},inner sep=2pt}
\tikzstyle{aut}=[shorten >=1pt,node distance=2cm,on grid,auto,initial text={},inner sep=2pt]
\tikzstyle{st}=[circle,draw=black,fill=black!10,inner sep=3pt]
\tikzstyle{sst}=[rectangle,draw=none,fill=none,inner sep=3pt]
\tikzstyle{final}=[accepting]

%%% Uppaal-like diagrams
\newcommand{\uppbox}[3][20mm]{\tikz{
  \node[black!15,fill=black!15,minimum width=#1,align=left](title){\textbf{{\footnotesize #2}}};
  \node[black!15,fill=black!15,left,xshift=4mm]at(title.east){\textbf{{\footnotesize #2}}};
  \node[blue!60!cyan,right] at(title.west){\textbf{{\footnotesize #2}}};
  \node[below,inner sep=2mm,fill=white,xshift=2mm](box)at(title.south){\includegraphics[width=#1]{#3}};
  \node[fit=(title)(box),draw=black,inner sep=0pt]{};
}}

\newcommand{\uppboxv}[3][20mm]{\tikz{
  \node[below,inner sep=2mm,fill=white](box){\includegraphics[height=#1]{#3}};
  \coordinate[yshift=5mm](top)at(box.north);
  \node[fit=(top)(box.north west)(box.north east),inner sep=0pt,fill=black!15](title){};
  \node[blue!60!cyan,right] at(title.west){\textbf{{\footnotesize #2}}};
  \node[fit=(title)(box),draw=black,inner sep=0pt]{};
}}


%% COnfiguring Listings
\lstset{ % basic style
  % language=scala,
  basicstyle=\ttfamily\scriptsize,
  breakatwhitespace=true,
  breaklines=true,
  mathescape,
  % morecomment=[l]{//},
  % morecomment=[n]{/*}{*/},
  % frame=single,                    % adds a frame around the code
  rulecolor=\color{black!40},         % if not set, the frame-color may be changed on line-breaks within not-black text (e.g. comments (green here))
  xleftmargin=1.5mm,
  xrightmargin=1.5mm,
  backgroundcolor=\color{black!5},
  % line numbers
%  numbers=left,  % where to put the line-numbers; possible values are (none, left, right)
 numbersep=5pt, % how far the line-numbers are from the code
 numberstyle=\tiny\color{gray},   
 stepnumber=1,  % the step between two line-numbers. If it is 1 each line will be numbered      
%  xleftmargin=3mm,
%  xrightmargin=1.5mm,
%%%%%
  captionpos=b, % t or b (top or bottom)
  belowcaptionskip=5mm,
%%%%%
  % alsoletter={-},
  % emphstyle=\ttfamily\color{blue}, %\underbar,
  % emphstyle={[2]\ttfamily\color{green!50!black}},
  emphstyle=\bfseries\itshape\color{blue!80!black},       % moreemph={...} - layer keywords
  emphstyle={[2]\itshape\color{red!70!black}},%\underbar} % moreemph={[2]...} - inner keywords
  %
  keywordstyle=\bf\ttfamily\color{red!50!black},
%  commentstyle=\sl\ttfamily\color{gray!70},
  commentstyle=\color{green!60!black},
  stringstyle=\ttfamily\color{purple!60!black},
  morestring=[b]",
  morecomment=[l]{\#},
  frame=single,
  % numberstyle=\tiny, numbers=left, stepnumber=1, firstnumber=1, numberfirstline=true,
  %emph={act,proc,init,sort,map,var,eqn},
  %emph={[2]block,hide,comm,rename,allow,||,<>,sum,&&,=>,true,false},
  % literate=*{->}{{{\color{red!70!black}$\to$}}}{1}
  %            {.}{{{\color{red!70!black}.}}}{1}
  %            {+}{{{\color{red!70!black}\hspace*{1pt}+\hspace*{1pt}}}}{1}
  %            {|}{{{\color{red!70!black}|}}}{1}
  %            {||}{{{\color{red!70!black}|\!\!|}}}{1}
  %            {*}{{{\color{red!70!black}*}}}{1}
  %            {\#}{{{\color{red!70!black}\#}}}{1}
  %            {&}{{{\color{red!70!black}\&}}}{1}
  %            {:=}{{{\color{red!70!black}:=}}}{1}
  %            {=>}{{{\color{red!70!black}=>}}}{1}
  %            % {>}{{{\color{green!65!black}\hspace*{1pt}>\hspace*{1.5pt}}}}{2}
  %            % {<}{{{\color{green!65!black}\hspace*{1.5pt}<\hspace*{1pt}}}}{2}
  %            % {]}{{{\color{green!65!black}\hspace*{1pt}]\hspace*{1.5pt}}}}{1}
  %            % {[}{{{\color{green!65!black}\hspace*{1.5pt}[\hspace*{1pt}}}}{1}
  %            {>}{{{\color{green!65!black}>}}}{1}
  %            {<}{{{\color{green!65!black}<}}}{1}
  %            {]}{{{\color{green!65!black}]}}}{1}
  %            {[}{{{\color{green!65!black}[}}}{1}
  %  morekeywords={Merger1,Fifo2,Lossy3,Init1,Init2}
  % ,emph={act,proc,init,sort,map,var,eqn}
  % ,emph={[2]block,hide,comm,rename,allow,||,<>,sum,&&,=>}
}

\lstdefinestyle{tiny}{basicstyle=\ttfamily\relsize{-7}}

% Include slides from others
\newcommand{\byothers}[3]{{
\begin{frame}{}~\mycbox{\Large slides by #1\\pages #2}\end{frame}{}
\setbeamercolor{background canvas}{bg=}
\includepdf[pages=#2]{../../others/#3}
}}

%-------------- template --------------------------------------------------
\usetheme{metropolis}
\usepackage{appendixnumberbeamer}

% Base colors (from metropolis theme)
\definecolor{metDarkBrown}{HTML}{604c38}
\definecolor{metDarkTeal}{HTML}{23373b}
\definecolor{metLightBrown}{HTML}{EB811B}
\definecolor{metLightGreen}{HTML}{14B03D}

 

\metroset{numbering=fraction,progressbar=frametitle}

% \setbeamercolor*{structure}{fg=blue!80!black}
\setbeamercolor*{structure}{fg=metLightGreen}

% % \definecolor{MainColour}{rgb}{0., 0.25, 0.8}
% \colorlet{MainColour}{blue!50!black}
% \colorlet{BgColour}{blue!10}
% \colorlet{BarColour}{blue!50!black}


% %\usetheme{CambridgeUS}%{Copenhagen}%{Frankfurt}%{Singapore}%{CambridgeUS}
% \usecolortheme[named=MainColour]{structure} 
% \useoutertheme[subsection=false]{miniframes}
% \useinnertheme{circles}
% %\useinnertheme[shadow=false]{rounded}
% \setbeamertemplate{blocks}[rounded][shadow=false]

% \setbeamercovered{transparent} 
% \setbeamertemplate{navigation symbols}{} %Remove navigation bar
% \setbeamertemplate{footline}[frame number] % add page number
% \setbeamercolor{postit}{fg=MainColour,bg=BgColour}
% \setbeamercolor{structure}{bg=black!10}
% %\setbeamercolor{palette primary}{use=structure,fg=red,bg=green}
% %\setbeamercolor{palette secondary}{use=structure,fg=red!75!black,bg=green}
% \setbeamercolor{palette tertiary}{use=structure,bg=BarColour,fg=white}
% %\setbeamercolor{palette quaternary}{fg=black,bg=green}
% %\setbeamercolor{normal text}{fg=black,bg=white}
% %\setbeamercolor{block title alerted}{fg=red,bg=green}
% %\setbeamercolor{block title example}{bg=black!10,fg=green}
\setbeamercolor{block body}{bg=black!5}

% \setbeamercolor{block title alerted}{bg=red!25}
% \setbeamercolor{block body alerted}{bg=red!10}

% \setbeamercolor{block title example}{bg={rgb:green,2;black,1;white,5}}
\setbeamercolor{block body example}{bg={rgb:green,2;black,1;white,20}}
\setbeamercolor{block body alerted}{bg={metLightBrown!25}}

% \setbeamertemplate{itemize item}{\color{black!10}$\blacksquare$}
\setbeamercolor{itemize item}{fg=metDarkTeal}
\setbeamercolor{itemize subitem}{fg=metDarkTeal}

\setbeamercolor{graybc}{fg=black,bg=black!10}
\newcommand{\myblock}[1]{\begin{beamercolorbox}[dp=1ex,center,rounded=true]%
  {graybc} {\large \textbf{#1}} \end{beamercolorbox}}%


% Configuring the foot line
\setbeamertemplate{footline}
{
  \leavevmode%
  \hbox{%
  \begin{beamercolorbox}[wd=.4\paperwidth,ht=2.25ex,dp=1ex,center]{author in head/foot}%
    \usebeamerfont{author in head/foot}\insertshortauthor
  \end{beamercolorbox}%
  \begin{beamercolorbox}[wd=.5\paperwidth,ht=2.25ex,dp=1ex,center]{title in head/foot}%
    \usebeamerfont{title in head/foot}\insertsection
  \end{beamercolorbox}%
  \begin{beamercolorbox}[wd=.1\paperwidth,ht=2.25ex,dp=1ex,right]{date in head/foot}%
    \insertframenumber{} / \inserttotalframenumber\hspace*{2ex} 
  \end{beamercolorbox}}%
  \vskip0pt%
}
% No configuration symbols
\setbeamertemplate{navigation symbols}{}


%----------------------------------------------------------------------------

\begin{document}

\setLecture{3}{Counting steps\\~~~(Asymptotic analysis)}
\frame[plain]{\titlepage}


\begin{frame}[t]\frametitle{Overview}

  % \begin{itemize}
  %   \item Measuring \structure{precisely} performance of algorithms
  %   \item Measuring \alert{asymptotically} performance of algorithms
  %   \item Measuring precisely the \alert{average time} of algorithms
  %   \item Next: beyond worst-/best-case scenarios
  %         \begin{itemize}
  %           \item Analysing \alert{recursive} functions
  %           \item Analysis of sequences of operations (\alert{amortised analysis})
  %         \end{itemize}
  % \end{itemize}

  \begin{itemize}
    {\color{black!25}
    \item Checking {correctness} of algorithms}
    \item Measuring \structure{precisely} performance of algorithms
    \item Measuring \alert{asymptotically} performance of algorithms
    \item Analysing \alert{recursive} functions
    \item Next: beyond worst-/best-case scenarios
          \begin{itemize}
            \item \alert{average time} of a single operation
            \item analysis of sequences of operations (\alert{amortised analysis})
          \end{itemize}
    \end{itemize}
\end{frame}

\section{Motivation}

\byothers{Charles E. Leiserson}{3-19}{Cormen_Algorithms_Slides.pdf}

\byothers{Pedro Ribeiro, slides 2}{1-2}{ribeiro/2_asymptoticanalysis_02102018.pdf}
\byothers{Pedro Ribeiro, slides 2}{8-18}{ribeiro/2_asymptoticanalysis_02102018.pdf}

\begin{frame}\frametitle{Next steps}
  \begin{enumerate}
    \item Precise analysis: \alert{counting operations}
    \item Approximate analysis -- \structure{Asymptotic notation }$(O, \Theta, \Omega, o, \omega)$
  \end{enumerate}
\end{frame}

\section{Counting operations}

% \begin{frame}[t]\frametitle{Exercises}
%   [WiP: bubble sort, iSort, mult1, mult2]

%   [Recall arithmetic and geometric series, height of binary tree, ...]

%   [Proceed with Pedro's slides]
% \end{frame}

\begin{frame}[fragile]\frametitle{Simpler counting}
\centering

\begin{lstlisting}[language=C++,linewidth=50mm]
int count = 0;
for (int i=0; i<n; i++)
  if (v[i] == 0) count++
\end{lstlisting}    

\frsplitt{
  \begin{block}{RAM}
  \begin{itemize}
    \item worst-case: $T(n) = 5+5n$
    \item best-case: $T(n) = 5+4n$
  \end{itemize}  
  \end{block}
}{
  \begin{alertblock}{\#array-accesses + \#count-increments}
  \begin{itemize}
    \item worst-case: $T(n) = 2n$
    \item best-case: $T(n) = n$
    \item \textcolor{black!30}{
      average-case: $\overline{T}(n) = n + \sum_{0\leq r < n}P(v[r]=0)$}
  \end{itemize}  
  \end{alertblock}
  
}
% 8kB
\end{frame}

\begin{frame}[fragile]\frametitle{Exercises}
  
\begin{columns}
\begin{column}{0.48\textwidth}
\begin{lstlisting}[language=C++,emph={bubbleSort,swap}]
void bubbleSort(int v[], int N){
  int i, j;
  for (i=N-1; i>0; i--)
    for (j=0; j<i; j++)
      if (v[j] > v[j+1])
        swap(v,j,j+1);
}
\end{lstlisting}
\end{column}
\begin{column}{0.48\textwidth}
\begin{lstlisting}[language=C++,emph={iSort,swap}]
void iSort(int v[], int N){
  int i, j;
  for (i=1; i<N; i++)
    for (j=i; j>0 && v[j-1]>v[j];
         j--)
      swap(v,j,j-1);
}
\end{lstlisting}
\end{column}
\end{columns}

\doSimpleExercise{What is the best and worst case wrt \textbf{comparisons between array elements}?\\[3mm]}
~\\[-12mm]
\doSimpleExercise{What is the best and worst case wrt \textbf{swaps}?}
~\\[-12mm]
\doSimpleExercise{\textbf{How many} of these operations are performed in both cases?}
\end{frame}

\begin{frame}[fragile]{Exercises}
  
\begin{columns}
\begin{column}{0.48\textwidth}
\begin{lstlisting}[language=C++,emph={mult1}]
int mult1 (int x, int y){
  int a, b, r;
  a=x; b=y; r=0;
  while (a!=0){
    r = r+b;
    a = a-1;
  }
  return r;
}
\end{lstlisting}
\end{column}
\begin{column}{0.48\textwidth}
\begin{lstlisting}[language=C++,emph={mult2}]
int mult2 (int x, int y){
  int a, b, r;
  a=x; b=y; r=0;
  while (a!=0) {
    if (a%2 == 1) r = r+b;
    a=a/2;
    b=b*2;
  return r;
}
\end{lstlisting}
\end{column}
\end{columns}
\doSimpleExercise{In each case, how many primitive operations (\texttt{+ - *2 /2 \%2}) are performed in the worst case?}
Note: In \texttt{mult2}, consider the size $N$ as the number of bits used to represent \texttt{x} and \texttt{y}; e.g., with 5 bits you can represent a positive integer until 31.

% Considere que o tamanho do input  ́e o nu ́mero de bits necess ́arios para representar os nu ́meros inteiros passados como argumento. Recorde que, por exemplo, os nu ́meros cuja representa ̧c ̃ao requer 5 bits s ̃ao {16, . . . , 31}.}
\end{frame}


\begin{frame}[fragile]\frametitle{Exercises}
\begin{columns}
\begin{column}{0.48\textwidth}
\begin{lstlisting}[language=C++,emph={maxgrow,grow}]
int maxgrow(int v[], int N) {
  int r = 1, i = 0, m;
  while (i<N-1) {
    m = grow(v+i, N-i);
    if (m>r) r = m;
    i++;
  }
  return r;
}
\end{lstlisting}
\end{column}
\begin{column}{0.48\textwidth}
\begin{lstlisting}[language=C++,emph={grow}]
int grow(int v[], int N) {
  int i;
  for (i=1; i<N; i++)
    if (v[i] < v[i-1]) break;
  return i;
}

\end{lstlisting}
\end{column}
\end{columns}

~\\[-6mm]
\doSimpleExercise{What is the best and worst case for \texttt{maxgrow} wrt comparisons of array elements?}
~\\[-11mm]
\doSimpleExercise{How many comparisons are in each case?}
~\\[-12mm]
\doSimpleExercise{If we replace \texttt{i++} by \texttt{i+=m}, how many comparisons are in the worst case?}

% Worst case when i+=m replaces i++?
\end{frame}


\begin{frame}[fragile]\frametitle{Exercises @home}
\begin{columns}
\begin{column}{0.48\textwidth}
\begin{lstlisting}[language=C++,emph={maxSum,sum}]
int maxSum(int v[], int N) {
  int i, j, r=0, m;
  for (i=0; i<N; i++)
    for (j=i; j<N; j++) {
      m = sum(v,i,j);
      if (m>r) r = m;
    }
  return r;
}
\end{lstlisting}
\end{column}
\begin{column}{0.48\textwidth}
\begin{lstlisting}[language=C++,emph={sum}]
int sum(int v[], int a, int b) {
  int r = 0, i;
  for (i=a; i<=b; i++)
    r=r+v[i];
  return r;
}


\end{lstlisting}
\end{column}
\end{columns}
%
% T(N) = ...
% =   sum_i=0^(N-1) (N-i)*((N-1-i)/2)
% =   1/2 * sum_i=0^(N-1) (N-i)(N-1-i)
% =   1/2 * sum_i=0^(N-1) (N^2-N-iN -iN +i +i^2)
% =   1/2 * (sum (N^2-N) - sum iN - sum iN + sum i + sum i^2)
% =   1/2 * (N(N^2-N) - N*sum i + sum i^2)
% = Theta(N^3)


% ~\\[-10mm]
\doSimpleExercise{What is the complexity \texttt{maxSum} wrt accesses to the array?}
\end{frame}


\section{Asymptotic Notation}

\byothers{Pedro Ribeiro, slides 2}{19-23}{ribeiro/2_asymptoticanalysis_02102018.pdf}

\begin{frame}[t]\frametitle{Examples}
    ~\\[-8mm]
    \begin{block}{Big Oh ($\mathcal{O}$)}
        $3n^2 -100n+6 \only<1>{\alert{~{=}?~}}\only<2->{\alert<2>{=   }} \mathcal{O}(n^2)$\only<2->{\structure<2>{ ~~because~~} $ 3n^2 >3n^2 -100n+6$}
      \\$3n^2 -100n+6 \only<1>{\alert{~{=}?~}}\only<2->{\alert<2>{=   }} \mathcal{O}(n^3)$\only<2->{\structure<2>{ ~~because~~} $0.01n^3 >3n^2 -100n+6$}
      \\$3n^2 -100n+6 \only<1>{\alert{~{=}?~}}\only<2->{\alert<2>{\neq}} \mathcal{O}(n)$  \only<2->{\structure<2>{~~~because~~} $c\cdot n<3n^2$ when $n>c$}
    \end{block}

    \begin{block}{Big Omega ($\Omega$)}
        $3n^2 -100n+6 \only<1-2>{\alert{~{=}?~}}\only<3->{\alert<3>{=   }} \Omega(n^2)$\only<3->{\structure<3>{ ~~because~~} $2.99n^2 <3n^2 -100n+6$}
      \\$3n^2 -100n+6 \only<1-2>{\alert{~{=}?~}}\only<3->{\alert<3>{\neq}} \Omega(n^3)$\only<3->{\structure<3>{ ~~because~~} $n^3 > 3n^2 -100n+6$}
      \\$3n^2 -100n+6 \only<1-2>{\alert{~{=}?~}}\only<3->{\alert<3>{=   }} \Omega(n)$  \only<3->{\structure<3>{~~~because~~} $10^{10^{10}} n < 3n^2 -100+6$}
    \end{block}

    \begin{block}{Big Theta ($\Theta$)}
        $3n^2 - 100n + 6 \only<1-3>{\alert{~{=}?~}}\only<4->{\alert<4>{=   }} \Theta(n^2)$\only<4->{\structure<4>{ ~~because~~} $\mathcal{O}$ and $\Omega$}
      \\$3n^2 - 100n + 6 \only<1-3>{\alert{~{=}?~}}\only<4->{\alert<4>{\neq}} \Theta(n^3)$\only<4->{\structure<4>{ ~~because~~} $\mathcal{O}$ only}
      \\$3n^2 - 100n + 6 \only<1-3>{\alert{~{=}?~}}\only<4->{\alert<4>{\neq}} \Theta(n)$  \only<4->{\structure<4>{~~~because~~} $\Omega$ only}
    \end{block}

\end{frame}

\byothers{Pedro Ribeiro, slides 2}{24-31}{ribeiro/2_asymptoticanalysis_02102018.pdf}


% \begin{frame}\frametitle{Exercises}

%   [WiP]

% \end{frame}

\byothers{Pedro Ribeiro, exercises 2}{1-2}{ribeiro/ex02.pdf}



\section{Analysis of recursive functions}

\begin{frame}[fragile]\frametitle{Binary search}

To analyse the complexity of a recursive function, we typically define the time $T$ using \emph{recurring equations}.

\begin{columns}
\begin{column}{0.48\textwidth}
\begin{lstlisting}[language=C++,emph={bsearch}]
int bsearch(int x, int v[], int N){
  int i;
  if (N<=0) i = -1;
  else {
    m = N/2;
    if (v[m]==x) i = m;
    else if (v[m] > x)
      i = bsearch(x, v, m);
    else {
      i = bsearch(x, v+m+1, N-m-1);
      if (i!=-1) i = i+m+1
    }
  return i ;
}
\end{lstlisting}
\end{column}
  \begin{column}{0.48\textwidth}
    \small
    \only<1>{Counting the number of comparisons with array elements:
    $$T(N) = \left\{\begin{array}{ll}
    0 & \text{if }N=0\\
    T(N/2)+2 & \text{if }N>0
    \end{array}\right.$$}
    \only<2->{
    (T(N) = T(N/2)+2~~if~~N>0)
    \begin{align*}
    T(0) &= 0\\
    T(1) &= T(2^0) = 2\\
    T(2) &= T(2^1) = 2+T({2}/{2}) = 2+2 = 4\\
    T(4) &= T(2^2) = 2+T({4}/{2}) = 2+2+2 = 6\\
    \cdots\\
    T(2^i) &= \underbrace{2+2+\cdots+2}_{i\text{-times}} + 2 = 2i+2
    \\
    \only<3>{T(N) &= T(2^{log_2(N)}) \\&= 2*log_2(N)+2 ~~~(= \Theta(log(N)))}
    \end{align*}
    }
  \end{column}
\end{columns}
\end{frame}

\byothers{Pedro Ribeiro, slides 2}{38-61}{ribeiro/2_asymptoticanalysis_02102018.pdf}


\begin{frame}[fragile]\frametitle{Exercises}

\doExercise{Solve using recursion trees (assume $T(0)$ is a constant)}{
  \small
  \begin{enumerate}
    \item $T(n)~=~k+T(n-1)$ ~~where $k$ is a constant
    \item $T(n)~=~k + T(n/2)$ ~~where $k$ is a constant
    \item $T(n)~=~k + 2 * T(n/2)$ ~~where $k$ is a constant
    \item $T(n)~=~n+T(n-1)$
    \item $T(n)~=~n + T(n/2)$
    \item $T(n)~=~n+2*T(n/2)$
  \end{enumerate}
}
\end{frame}

\begin{frame}[fragile]\frametitle{Exercises}
~\\[-8mm]
\doSimpleExercise{Write recurrences for \texttt{maxSumR} (wrt array accesses) and \texttt{Hanoi} (wrt \texttt{printf})}
~\\[-12mm]
\doSimpleExercise{Draw a recurrence tree for \texttt{Hanoi} and use it to derive its asymptotic complexity}
~\\[-8mm]
\begin{columns}
\begin{column}{0.48\textwidth}
\begin{lstlisting}[language=C++,emph={maxSumR}]
int maxSumR (int v[], int N) {
  int r=0, m1, m2, i;
  if (N>0) {
    m1 = m2 = v[0];
    for (i=1; i<N; i++) {
      m2 = m2+v[i];
      if (m2>m1) m1=m2;
    }
    m2 = maxSumR (v+1,N-1);
    if (m1>m2) r = m1; else r = m2;
  }
  return r; }
\end{lstlisting}
% T(0) = 0 | T(n-1) + n-1
\end{column}
\begin{column}{0.48\textwidth}
\begin{lstlisting}[language=C++,emph={Hanoi}]
void Hanoi(int nDiscs, int l,
           int r, int m) {
  if (nDiscs > 0) {
    Hanoi(nDiscs-1, l, m, r);
    printf("move disk: %d --> %d\n", l, r);
    Hanoi(nDiscs-1, m, r, l);
  }
}
\end{lstlisting}
\end{column}
\end{columns}

\end{frame}

\begin{frame}[fragile]\frametitle{Exercises}
\begin{columns}
\begin{column}{0.51\textwidth}
\begin{lstlisting}[language=C++,emph={heightBT}]
int heightBT(BTree t){
  int r=0;
  if (t!=NULL)
    r = 1 + max (heightBT(t->left),
                 heightBT(t->right));
  return r;
}
\end{lstlisting}
\end{column}
\begin{column}{0.45\textwidth}
\doSimpleExercise{Recall binary trees; this function calculates the maximum height of a binary tree. Identify the best and worst cases for this function, and describe a recurrence for each one.}
\end{column}
\end{columns}
\end{frame}

% T(N) = 
% k1 if t==null
% k2 + 2T(N/2)   (best case)
% T(N) =
% k1 if t==null
% k2 + T(N-1)    (worst case)
%%% use "comparisons with null": 1 | 1+2T(N/2)

% \begin{frame}\frametitle{Exercises}
% [more exercises: maxSumR; hanoi; heightBT]
% [maybe use PR's exercises too.]
% \end{frame}


\byothers{Pedro Ribeiro, slides 2}{62-69}{ribeiro/2_asymptoticanalysis_02102018.pdf}

\byothers{Pedro Ribeiro, exercises 3}{1-2}{ribeiro/ex03.pdf}

% \begin{frame}\frametitle{Exercises}
%   \doSimpleExercise{Calculate the asymptotic complexity of the 3 exercises above (\texttt{maxSumR}, \texttt{hanoi}, and \texttt{heightBT}), indicating which case you used.}
% \end{frame}

% \begin{frame}[t]\frametitle{...}
    
% [example with merge sort; avoid substitution method (or glimpse for a given guess); focus on recurrent trees]

% [exercises both from Ribeiro and JBB]

% [Later: simplified version of master theorem and short exercises]

% \end{frame}




\end{document}
